% Generated draft: requires mathematical and source review.
\documentclass[11pt]{article}
\usepackage[margin=25mm]{geometry}
\usepackage{fontspec}
\setmainfont{Latin Modern Roman}
\usepackage{amsmath,amssymb,mathtools,unicode-math}
\setmathfont{Latin Modern Math}
\usepackage{enumitem,needspace,xurl,hyperref}
\setlength{\parindent}{0pt}
\setlength{\parskip}{5pt}
\setlength{\emergencystretch}{3em}
\newcommand{\N}{\mathbb{N}}
\newcommand{\Z}{\mathbb{Z}}
\newcommand{\Q}{\mathbb{Q}}
\newcommand{\R}{\mathbb{R}}
\newcommand{\C}{\mathbb{C}}
\newcommand{\F}{\mathbb{F}}
\newcommand{\A}{\mathbb{A}}
\newcommand{\PP}{\mathbb{P}}
\newcommand{\E}{\mathbb{E}}
\newcommand{\eps}{\varepsilon}
\newcommand{\dd}{\,\mathrm{d}}
\begin{document}
\section*{UnsolvedMath ID: 30006296}
\textbf{Source problem code:} catalogue entry
\subsection{Definitions and mathematical statement}
For terminology, a finite matroid is a pair \(M=(E,\mathcal I)\), where \(E\) is a finite set and \(\mathcal I\subseteq 2^E\) satisfies \(\varnothing\in\mathcal I\), closure under taking subsets, and the augmentation axiom: if \(I,J\in\mathcal I\) and \(|I|<|J|\), then some \(e\in J\setminus I\) satisfies \(I\cup\{e\}\in\mathcal I\). This defines the finite case for context; the supplied record does not explicitly specify the intended class of matroids or impose representability or looplessness.

A smooth projective toric variety over a field \(k\) is a smooth projective variety \(X\) containing a split algebraic torus \(T\cong(\mathbb G_{m,k})^d\) as a dense open subset, such that the multiplication action of \(T\) on itself extends to an action on \(X\). The supplied record does not specify a base field or a convention concerning split versus nonsplit toric varieties.

\textbf{Research question.} Develop a precise theory describing the extent to which matroids admit structures analogous to those of smooth projective toric varieties, including analogues of sheaf cohomology, vector bundles, intersection theory, and Hodge-theoretic structures.

This is a research program rather than a specified conjectural proposition. The source supplies neither definitions of the proposed matroid-side sheaves and vector bundles nor comparison constructions, compatibility requirements, or quantified assertions identifying what such a theory must establish. In particular, it does not assert that every matroid is realized by a smooth projective toric variety. The intended class of matroids and the precise comparison targets remain unresolved.

The supplied dated assessment reports strong toric-like Hodge properties for matroid Chow rings, while describing the broader program as unresolved. It gives no precise Hodge statements or hypotheses; this contextual assessment does not specify or settle the requested analogues.
\subsection{Short English statement}
The question asks how far the geometry of smooth projective toric varieties can be reproduced for matroids, particularly for sheaf cohomology, vector bundles, intersection theory, and Hodge theory. The supplied formulation is a broad research program: its objects, hypotheses, and required comparison properties must be specified before it becomes a precise conjecture. Its open status is supplied curation, not an independently verified conclusion.
\subsection{Sources}
\begin{itemize}
\item Local UnsolvedMath registry, supplied complete source record for permanent ID \(30006296\), title \emph{Matroids as Smooth Projective Toric Varieties}, source problem code \texttt{catalogue entry}; collection: Oberwolfach Reports Open Problems. The record lists proposed year 2025 and no named proposer.
\item \emph{Toric Geometry} (2025), Oberwolfach Reports, pp.\ 883--940, DOI \texttt{10.4171/owr/2025/19}: \url{https://doi.org/10.4171/owr/2025/19}. This is the supplied report provenance; no exact problem page or full report passage is included in the supplied record.
\item Supplied literature source, labelled ``Eur, matroids and projective toric varieties'': \url{https://ems.press/content/serial-article-files/51856}.
\item Supplied literature source, labelled ``Eur, conference abstract'': \url{https://lsa.umich.edu/math/news-events/all-events.detail.html/125630-21855469.html}.
\item Literature assessment dated \texttt{2026-08-22}: supplied dated curation, not independently verified. It labels the problem open, reports strong toric-like Hodge properties of matroid Chow rings, and describes the broader question concerning sheaf cohomology, vector bundles, and a full toric analogue as unresolved and not a single settled claim. The source explicitly requests independent verification of time-sensitive status and cited claims.
\end{itemize}
\end{document}
